\chapter{Despliegue y reproducibilidad}

Desde el primer sprint, \texttt{docker compose up --build} crea tres servicios. PostgreSQL mantiene un volumen persistente y expone una comprobación de salud; la API espera a que la base esté disponible, aplica migraciones y escucha en el puerto 8080; Nginx sirve la SPA en el puerto 5173 y reenvía \texttt{/api} al backend.

Los valores del Compose permiten una demostración local inmediata. El fichero \texttt{.env.example} enumera las variables que deben sustituirse y \texttt{.env} queda excluido del control de versiones. En particular, claves JWT, token de bootstrap y contraseña de base de datos deberán gestionarse como secretos en cualquier entorno accesible.

El objetivo de este sprint es reproducibilidad local, no operación pública. La selección del destino accesible, HTTPS, copias de seguridad, observabilidad mínima y automatización de despliegue pertenece al Sprint 7.
